% TOP_PROBLEM: {"problemId": "problem.algebraic-independence-of-e-and-pi", "canonicalTitle": "Algebraic Independence of e and pi", "releaseRank": 168, "primaryDomain": "number_theory_arithmetic_geometry", "primaryDomainLabel": "Number theory & arithmetic geometry", "displayStatus": "open", "releaseStatus": "open"}
% !TeX program = lualatex
% Definition and sources only; this file is not an LLM solution attempt.
\documentclass[11pt]{article}
\usepackage[margin=25mm]{geometry}
\usepackage{fontspec}
\setmainfont{Latin Modern Roman}
\usepackage{amsmath,amssymb,mathtools}
\usepackage{unicode-math}
\setmathfont{Latin Modern Math}
\usepackage{xurl,hyperref}
\hypersetup{colorlinks=true,urlcolor=blue}
\setcounter{secnumdepth}{0}
\setlength{\parindent}{0pt}
\setlength{\parskip}{5pt}
\setlength{\emergencystretch}{3em}
\begin{document}
\section*{\texorpdfstring{Algebraic Independence of e and pi}{Algebraic Independence of e and pi}}\label{168}
\subsection{Definitions and mathematical statement}
Define $e=\sum_{n=0}^{\infty}1/n!$ and let $\pi$ be the positive half-period of the usual sine function. Algebraic independence of $e$ and $\pi$ over ${\mathbb{Q}}$ means
\[
 \forall P\in{\mathbb{Q}}[X,Y]\setminus\{0\},\qquad P(e,\pi)\ne0.
\]
Equivalently, the evaluation homomorphism ${\mathbb{Q}}[X,Y]\to{\mathbb{R}}$ is injective, or the transcendence degree of ${\mathbb{Q}}(e,\pi)$ over ${\mathbb{Q}}$ is two. This is stronger than showing either number individually transcendental or showing just $e+\pi$ irrational. A counterexample would require a nonzero polynomial relation with rational coefficients, not a numerical near-equality.
\par\smallskip\textit{Formulation and concept references:} \href{https://people.maths.ox.ac.uk/pila/LMSNotes.pdf}{[S1]}.

\subsection{Short English statement}
Are the constants e and pi free of every nontrivial polynomial relation with rational coefficients?
\subsection{Sources}
\begin{itemize}
\item[S1]Jonathan Pila, Functional transcendence via o-minimality, edited LMS-EPSRC 2013 lecture notes.
\url{https://people.maths.ox.ac.uk/pila/LMSNotes.pdf}.
\textit{Formulation source.}
\end{itemize}
\end{document}
